\section{Preliminaries}
\label{sec:prelim}

\subsection{Feed-forward pointmaps and Gaussians}
DUSt3R~\cite{wang2024dust3r} maps an uncalibrated image pair
$I^i, I^j \in \mathbb{R}^{H\times W\times 3}$ to pointmaps
$X^i_i, X^j_i \in \mathbb{R}^{H\times W\times 3}$ with confidences
$C^i_i, C^j_i$, where $X^i_j$ denotes the pointmap of image $i$ expressed in
the frame of camera $j$. MASt3R~\cite{leroy2024mast3r} adds a matching head
producing features $D^i_i, D^j_i \in \mathbb{R}^{H\times W\times d}$. We write
$F_M(I^i, I^j)$ for one forward pass.

Splatt3R~\cite{smart2024splatt3r} adds a third head, in parallel to the
pointmap and matching heads and using the same DPT decoder, which emits for
every pixel a rotation quaternion $q\in\mathbb{R}^4$, a scale
$s\in\mathbb{R}^3$, spherical-harmonic coefficients
$S\in\mathbb{R}^{3\times d}$, an opacity $\alpha\in\mathbb{R}$ and a mean
offset $\Delta\in\mathbb{R}^3$, so that the Gaussian centre is
\begin{equation}
\mu \;=\; x + \Delta, \qquad x \in X^i_i .
\label{eq:mu}
\end{equation}
Quaternions are normalised, scales and offsets pass through an exponential and
opacities through a sigmoid. Critically for this paper, \emph{only the Gaussian
head is trained}: the encoder and decoder remain those of the pre-trained
MASt3R model, and all Gaussian parameters are predicted in the first image's
camera frame, avoiding any ground-truth transformation at inference.

The covariance follows the standard 3DGS~\cite{kerbl20233dgs} parameterisation
\begin{equation}
\Sigma \;=\; R(q)\,\mathrm{diag}(s)^2\,R(q)^\top ,
\label{eq:cov}
\end{equation}
and a view with world-to-camera $W$ and projective Jacobian $J$ renders with
the projected covariance $\Sigma' = J W \Sigma W^\top J^\top$, compositing
front-to-back:
\begin{equation}
c \;=\; \sum_{k} c_k\,\alpha_k \prod_{l<k}\bigl(1-\alpha_l\bigr).
\label{eq:alphacomp}
\end{equation}
\Cref{eq:alphacomp} is the reason opacity behaves as it does in our analysis
(\cref{sec:method:mech}): with a black background, the residual transmittance
$\prod_l (1-\alpha_l)$ contributes zero radiance, so lowering $\alpha$ darkens
a region and raising it brightens it.

Splatt3R is supervised only by novel-view renderings, under a mask $M$ of
pixels visible in at least one context view:
\begin{equation}
\mathcal{L} = \lambda_{\text{MSE}}\,\mathcal{L}_{\text{MSE}}(M\!\odot\!\hat{I}, M\!\odot\!I)
            + \lambda_{\text{LP}}\,\mathcal{L}_{\text{LPIPS}}(M\!\odot\!\hat{I}, M\!\odot\!I),
\label{eq:splatt3rloss}
\end{equation}
with $\lambda_{\text{LP}}{=}0.25$ in the released configuration. Two properties
of \cref{eq:splatt3rloss} drive our method. It is \emph{purely photometric} ---
depth enters only through the mask $M$ and through pair selection, never
through a gradient on geometry --- and it is \emph{pairwise}: the loss never
observes more than two context views, hence never observes a fused map.

In the released implementation the DC colour is anchored to the source image:
the head predicts a residual which is added to the spherical-harmonic encoding
of the input RGB value,
\begin{equation}
S_{\mathrm{DC}}(u)=S_{\mathrm{res}}(u)+\operatorname{RGB2SH}(I^i(u)).
\label{eq:colorresidual}
\end{equation}
Consequently, revisits can inherit exposure, white-balance, shadow, and
view-dependent differences from their source frames. Exposure normalisation
removes only the global component of this source-frame photometric gauge; it
does not enforce equality of $S_{\mathrm{DC}}$ for the same world surface.

\subsection{SLAM back-end}
We inherit MASt3R-SLAM's~\cite{murai2025mast3rslam} front- and back-end
unchanged. Poses are $\mathrm{Sim}(3)$ to absorb the scale inconsistency of
network predictions,
\begin{equation}
T=\begin{bmatrix} sR & t\\ 0 & 1\end{bmatrix},\qquad
T \leftarrow \tau \oplus T \triangleq \mathrm{Exp}(\tau)\circ T ,
\label{eq:sim3}
\end{equation}
with $\tau\in\mathfrak{sim}(3)$. Under a generic central-camera assumption a
pointmap induces its own camera model through the ray map
$\psi(X^i_i)$ (unit-norm normalisation), and tracking minimises a ray error
rather than a 3D point error, because
$\lVert\psi_1-\psi_2\rVert^2 = 2(1-\cos\theta)$ is bounded and therefore
robust to the depth errors that pointmap prediction frequently exhibits. The
canonical keyframe pointmap is maintained by a confidence-weighted running
average,
\begin{equation}
\tilde X^k_k \leftarrow
\frac{\tilde C^k_k \tilde X^k_k + C^f_k\, T_{kf} X^f_k}{\tilde C^k_k + C^f_k},
\qquad
\tilde C^k_k \leftarrow \tilde C^k_k + C^f_k .
\label{eq:fusion}
\end{equation}
Loop closures enter a $\mathrm{Sim}(3)$ pose graph solved by Gauss--Newton.
We modify none of this; the consequence is that our tracking is identical to
MASt3R-SLAM by construction, which \cref{sec:exp:ate} verifies externally.

\section{Method}
\label{sec:method}

\subsection{System}
We replace the pointmap prediction with Splatt3R, so each keyframe pair yields
a per-pixel Gaussian in one forward pass. Keyframe Gaussians are stored in
their anchor keyframe's local frame and composed into the world through that
keyframe's \emph{current} pose,
\begin{equation}
\mu^{W} = A_k\,\mu^{k} + b_k, \qquad
\Sigma^{W} = A_k\,\Sigma^{k}\,A_k^{\top},
\label{eq:anchor}
\end{equation}
where $[A_k\,|\,b_k]$ is the $\mathrm{Sim}(3)$ matrix of keyframe $k$. Because
the composition reads the pose at render time, a loop closure that updates
$T_{WC_k}$ moves the map with it at no cost and without re-baking --- the
property our refinement design depends on
(\cref{sec:method:refiner}).

The implementation stores means, scales, rotations, SH coefficients, opacity,
confidence, and validity in a shared keyframe buffer. Tracker updates may
rewrite pointmap state without clearing the Gaussian record unless a new
Gaussian prediction is present. This invariant prevents a geometric update
from silently erasing an already renderable keyframe. Filtering and the
global Gaussian limit are applied before rasterisation; the limit is therefore
a rendering budget implemented through spatial stride, rather than FIFO
deletion of individual keyframes.

\subsection{Where adaptation is admissible}
\label{sec:method:head}

Splatt3R trains its head against \emph{frozen} encoder features, making the
head a function of one specific feature distribution. We measured what happens
when that distribution moves:

\begin{itemize}\itemsep2pt
\item \textbf{Encoder LoRA}~\cite{hu2022lora} ($r{=}8$, $\alpha{=}16$ on
      \texttt{qkv}/\texttt{proj}/\texttt{fc1}/\texttt{fc2}, $46.6$M trainable):
      $-49\%$ reconstruction quality against the released checkpoint, with
      predicted scales exploding by $1821\times$. The failure is immediate,
      not a late-training instability.
\item \textbf{Decoder-only LoRA}: $+0.37$\dB{} at epoch~1, then collapse to
      $-1.61$\dB{} by epoch~5 with a $12\times$ climb in the $99$th-percentile
      scale.
\item \textbf{Head-only} (encoder and decoder bit-frozen, only
      \texttt{gaussian\_dpt} trainable): improves, and is what we keep.
\end{itemize}

The interpretation is structural rather than a matter of hyper-parameters. The
scale head is exponentially parameterised, $s = \exp(\hat{s})$, so it is
unconstrained above; feeding the head out-of-distribution features perturbs
$\hat{s}$ in a regime where the photometric loss of \cref{eq:splatt3rloss} is
nearly flat --- a large, nearly transparent Gaussian and a small opaque one
render almost identically at the supervision views. Geometry is therefore free
to diverge while the loss barely moves, which is exactly the observed
signature (quality collapse with scale explosion). Head-only training keeps
$46.6$M parameters frozen and trains the $158$M-parameter head with the
upstream objective and Adam~\cite{kingma2015adam}, per dataset family.

\subsection{Injection-time thinning}
\label{sec:method:thin}

A SLAM map covers each surface with predictions from many keyframes. Under
\cref{eq:alphacomp} with a black background, a surface covered by many
over-solid Gaussians composites to haze rather than to a sharper surface. We
therefore fade opacity at the moment a Gaussian enters the map, in proportion
to how little the pointmap confidence vouches for it.

Confidences are not calibrated across keyframes, so we rank-normalise
\emph{within} each keyframe. For a keyframe with $N$ Gaussians and confidence
values $\{c_n\}$, let $r_n$ be the rank of $c_n$ and
$\hat{c}_n = r_n/(N-1) \in [0,1]$. The injected opacity is
\begin{equation}
\alpha'_n \;=\; \alpha_n \cdot
\mathrm{clip}\Bigl(1 - 2\lambda\,(1-\hat{c}_n),\; 0.1,\; 1\Bigr),
\label{eq:thin}
\end{equation}
with dose $\lambda{=}0.45$ shipped. The lower clip at $0.1$ is deliberate: it
keeps \cref{eq:thin} a \emph{thinning} prior rather than a deletion prior, so
that even the least-confident Gaussian retains gradient and the refiner can
recover it if the photometric evidence supports it.

\paragraph{Why a training-time penalty cannot substitute.}
An opacity penalty added to \cref{eq:splatt3rloss} would be a different
intervention, because crowding \emph{does not exist} at training time: it is a
property of the assembled map, which forms only after SLAM accumulates
keyframes over a surface, and \cref{eq:splatt3rloss} never sees more than two
context views. Consistent with this, we measured that the refiner
un-saturates opacity on its own even with an untrained head (median
$1.0\rightarrow0.26$), identifying thinning as an \emph{acceleration along a
path the optimiser already walks} --- which also explains its bounded
downside, since a prior the optimiser would have reached anyway cannot do much
harm. A dose--response check at two thinning depths (opacity knees $0.9$ and
$0.6$) gave $-21.0\%$ and $-15.8\%$ LPIPS; neither collapsed toward zero, as a
training-substitutable effect would require.

\subsection{Online map refinement}
\label{sec:method:refiner}

Prediction gives a good initialisation, not a converged map. We add a third
process that photometrically refines the assembled map \emph{while SLAM runs}.

\paragraph{Parameterisation.}
Gaussians are optimised in their anchor keyframe's local frame in
pre-activation form --- $\log s$, unnormalised $q$, $\mathrm{logit}\,\alpha$,
and the SH DC term --- so that scales stay positive and opacity stays in
$(0,1)$ without constraints, and composed to the world by
\cref{eq:anchor} at every step. Per-group Adam rates follow the 3DGS
reference, with the positional rate scaled by scene extent:
$\eta_\mu = 1.6\!\times\!10^{-4}\!\cdot\!\mathrm{extent}$,
$\eta_{f_{dc}} = 2.5\!\times\!10^{-3}$, $\eta_\alpha = 5\!\times\!10^{-2}$,
$\eta_s = 5\!\times\!10^{-3}$, $\eta_q = 1\!\times\!10^{-3}$.

\paragraph{Objective.}
Supervision views are tracked frames, so the refiner optimises
\begin{equation}
\mathcal{L}_{\text{ref}} = (1-\lambda_{D})\,\lVert \hat{I}-I \rVert_1
                          + \lambda_{D}\,\bigl(1-\mathrm{SSIM}(\hat{I},I)\bigr),
\qquad \lambda_{D}=0.2 ,
\label{eq:refloss}
\end{equation}
i.e.\ the 3DGS objective rather than \cref{eq:splatt3rloss}: LPIPS is a
reporting metric here, not a training signal, so that our evaluation metric is
not also our loss.

\paragraph{Anchor-carried supervision.}
Each supervision frame is stored as an image together with
$(\text{anchor index}, T_{\text{anchor},f})$ rather than a world pose. Its
world pose is recomposed through the anchor's \emph{current} pose at sample
time, so supervision follows loop closures exactly as the map does under
\cref{eq:anchor}. Frames are held in CPU shared memory, which is what allows
the refiner to run on a second GPU without pinning both processes to one card.
If frame $f$ is anchored to keyframe $k$, the stored relative transform gives
\begin{equation}
T_{WC_f}=T_{WC_k}T_{C_kC_f}.
\label{eq:supervisionpose}
\end{equation}
The refiner consequently never becomes a second pose-graph writer: a loop
closure changes the anchor pose once, and both map placement and supervision
views follow that correction.

\paragraph{Sampling.}
Supervision is drawn from a reservoir of $200$ frames spanning the whole
sequence plus a ring of the $64$ most recent, with a recent fraction of $0.3$.
We measured uniform-over-history to dominate a recency-weighted schedule
(held-out early/late split: $14.40/14.09/14.72$ uniform versus
$14.06/13.79/14.35$ recent-only), because causal arrival already biases the
pool toward recent views.

\paragraph{Scheduling.}
Sharing one GPU, the refiner holds a duty cycle $d$ by sleeping
$\mathrm{EMA}[\Delta t]\,(1-d)/d$ after each step, so optimisation steps are
dropped but never frames. Refiner load on the \emph{same} card doubles tracker
latency ($101\!\rightarrow\!206$\,ms p50); on a second card it is noise
($103$\,ms), which is the configuration we deploy. After the sequence ends an
optional polish phase continues with the final poses.

\paragraph{Frozen centres.}
We hold $\mu$ fixed during refinement. The centres come from the network's
pointmap, and photometric gradients degrade them: because the loss of
\cref{eq:refloss} is evaluated at the supervision views themselves, moving a
centre can reduce training error while worsening the geometry seen from new
viewpoints. This was measured optimal at both online ($\sim$300-step) and
offline (3000-step) budgets.

The refiner publishes a versioned CPU snapshot rather than mutable GPU state.
The visualiser uploads only complete versions and retains the previous valid
version on a read/write collision. If the display buffer is smaller than the
refined map, publication subsamples after the full map has been saved. This
separates bounded display memory from the evaluation artifact and prevents a
large-map publication failure from discarding a valid refinement result.

\subsection{What opacity is}
\label{sec:method:mech}

The two preceding levers are, we argue, one phenomenon. Consider what the
architecture permits. The DC colour is anchored to the input pixel plus a
small learned residual and is sigmoid-bounded, so it cannot brighten a region
much past what its input pixel suggests. Under \cref{eq:alphacomp} with a black
background, $\alpha$ can: lowering it darkens toward black, raising it
brightens by admitting less background. Opacity is therefore the only channel
with headroom to \emph{lift} content --- a brightness-trim dial, not the
confidence or hedging gate it is usually taken to be.

This makes a falsifiable, sign-specific prediction. Let $\ell$ be input
luminance and $\rho$ the Spearman correlation between $\alpha$ and $\ell$ over
predicted Gaussians. A head trained with independently darkened inputs
(synthetic vignetting and white-balance jitter, where the target view's
darkening cannot be inferred from the input) must learn to lift what the input
darkened, so $\rho<0$. A head trained on real sensor data with no synthetic
darkening should instead place solid Gaussians on bright, reliably
reconstructable content, so $\rho>0$. We registered both signs before
measuring; \cref{sec:exp:mech} reports the result.
